\documentclass[../main.tex]{subfiles}
\begin{document}
%==========================================TIÊU ĐỀ TẬP========================================
\begin{table}[h]
    \begin{adjustwidth}{-1cm}{}
        \begin{tabular}{>{\centering\arraybackslash}p{6cm}|>{\raggedright\arraybackslash}p{12.5cm}}
            \multirow{5}{6cm}
            {\\[-40pt]
                \flaregothic\fontsize{20pt}{20pt}\selectfont \centering
                \color{eptype}HISTOIRE\color{black}\\
                \freshbot\fontsize{72pt}{72pt}\selectfont
                \color{epnum}24\\[6pt]
                \cafeteria\fontsize{20pt}{20pt}\selectfont\color{epnum}(D-24)\color{black}
                \\[18pt]
                \flaregothic\fontsize{18pt}{18pt}\selectfont
                \color{epattr}RÉUNION
                \color{black}
            }
             &
            {\fotskip\fontsize{18pt}{18pt}\selectfont \ruby{永}{えい}\ruby{遠}{えん}の\ruby{友}{ゆう}\ruby{情}{じょう}【holoFive】}
            \\[6pt]
             & {\cafeteria\fontsize{18pt}{18pt}\selectfont Forever Five!}                                                                                                                         \\[4pt]
             & {\vnmsans\fontsize{18pt}{18pt}\selectfont Tình bạn vĩnh cửu I: holoFive}                                                                                                           \\[8pt]
             & {\caviar{\fontsize{14pt}{14pt}\selectfont Nhân vật: \emp{delu} \emp{nene} \emp{lamy} \emp{botan} \emp{polka} \emp{kuon2} \emp{ryuuhou2} \emp{suzuno2} \emp{naomi2} \emp{kaigou2}}} \\ & (Phần truyện này được mượn phần lớn ý tưởng từ \href{https://archiveofourown.org/works/47009050}{Forever 5} trên AO3, viết bởi \href{https://archiveofourown.org/users/Estnov/pseuds/Estnov}{Estnov.})
            \\[8pt]
        \end{tabular}
    \end{adjustwidth}
\end{table}
%=========================================TRUYỆN CHÍNH========================================
\color{black}

\txtbx[2]{
    {\color{vol} \uline{Tóm tắt:}} Ký ức về ngày chia tay cuối tháng Tám năm 2020 vẫn là một vết cắt sâu trong lòng các thành viên Thế hệ thứ Năm và Tổng quản lý Kuon. Sau gần hai năm bền bỉ bước tiếp dưới bản thể Delutaya, cô gái nhỏ năm nào nay đã tìm lại niềm tin và đam mê âm nhạc. Một chiều cuối tháng Ba năm 2022 tại Kawachi, một cú vấp ngã bất cẩn của Nene trên đường mua đồ lẩu đã mở ra cuộc hội ngộ định mệnh. Từng lớp manh mối quen thuộc - từ chất giọng đặc trưng, chiếc móc khóa sừng quỷ đến những thói quen cử chỉ bất biến - đã kéo các mảnh ghép NePoLaBo lại gần nhau trong nước mắt và nụ cười. Giữa bàn lẩu shabu-shabu ấm áp của nhà Hikawa, lời thề về một holoFive trọn vẹn năm người mãi mãi được khắc sâu.
}

\pov{Hikawa Kuon}

Có những buổi chiều mà sự im lặng còn nặng nề hơn bất kỳ tiếng gào thét nào.

Tôi nhớ rất rõ cảm giác ấy vào ngày cuối cùng của tháng Tám năm 2020. Trong phòng thu của trụ sở công ty, ánh đèn trần trắng lạnh phản chiếu lên mặt bàn điều khiển âm thanh vắng bóng người. Chiếc ghế xoay bọc đệm ở góc phòng - nơi đáng lẽ phải có một cô gái nhỏ ngồi đung đưa đôi chân trong giờ nghỉ giải lao - để trống từ nhiều ngày trước. Chiếc tai nghe kiểm âm màu tím nhạt vẫn nằm ngay ngắn trên giá đỡ kim loại, không một vết bụi, nhưng cũng sẽ chẳng còn bàn tay nào với tới để đeo nó lên nữa.

Bản thông cáo nội bộ chỉ vỏn vẹn vài dòng ngắn ngủi. Không khí ngột ngạt bao trùm cả tầng hành lang sau khi cuộc họp báo khép lại. Những người đi qua nhau chỉ dám trao nhau ánh nhìn tránh né, khẽ cúi đầu thở dài. Mười sáu ngày hoạt động ít ỏi tựa như một cái chớp mắt giữa dòng chảy vội vã của giới giải trí. Giấc mơ vừa mới hé nở đã bị dập tắt không thương tiếc trước những đợt sóng ngầm vô hình từ sự tàn nhẫn của mạng xã hội.

Lúc ấy, mọi lời động viên hay an ủi đều trở nên vô nghĩa.

Nhưng cuộc sống thì không dừng lại để đợi một vết thương lành miệng. Khi cánh cửa hào nhoáng của sân khấu lớn khép lại trước mắt cô bé, tôi đã mở cánh cửa căn nhà mình tại Kawachi để đón em về. Em sợ phải quay về Quỷ giới trong sự tủi hổ, sợ ánh mắt phán xét của người đời. Căn phòng nhỏ ở tầng hai trở thành một chốn trú ẩn thầm lặng.

Nơi đó, theo năm tháng, không còn bóng tối của sự suy sụp. Tôi nghe thấy tiếng bàn phím lách cách gõ từng dòng beat nhạc, nghe nhịp chân nhè nhẹ trên sàn gỗ khi em tập lấy hơi, và rồi cả tiếng hát trong trẻo, mộc mạc lại bắt đầu cất lên sau tấm vách cách âm. Dưới cái tên Delutaya, em bền bỉ nhặt lại từng mảnh vỡ của niềm tin, tự bước đi bằng chính đôi chân mình như một nghệ sĩ độc lập. Em không còn là nàng tiểu quỷ với đôi cánh nhỏ mang thương hiệu của công ty, mà là một con người bằng xương bằng thịt đang sống kiên cường với đam mê âm nhạc.

Em đã tự chữa lành và bước tiếp. Nhưng còn bốn người đồng đội ở lại thì sao?

Ký ức của tôi bất giác quay ngược về buổi sáng ngày định mệnh ấy, trong căn phòng stream của Botan.

\ct

Ba mươi mốt tháng Tám năm 2020.

Bốn người họ ngồi thành một hàng trước màn hình máy tính. Không có mô hình 2D sống động chuyển động theo từng nụ cười, chỉ có một khung hình tĩnh với dòng chữ thông báo đơn sơ cùng biểu tượng của các thành viên Thế hệ thứ Năm.

Tôi đứng phía sau góc máy, tay cầm xấp khăn giấy trắng và chai nước lọc vừa mở nắp, lặng lẽ quan sát từng người.

Lamy ngồi ở chiếc ghế bên mép trái. Đôi bàn tay nhỏ nhắn của nàng Yêu tinh tuyết siết chặt lấy tà váy đến mức khớp ngón tay trắng bệch ra. Đôi mắt em đã hoe đỏ từ trước khi tín hiệu phát trực tiếp bắt đầu. Em cắn chặt môi dưới, cố ngăn những tiếng nấc nghẹn ngào chực trào lên cổ họng, nhưng từng giọt nước mắt vẫn không kìm được mà nối đuôi nhau lăn dài trên gò má.

Bên cạnh em, Nene cúi gằm mặt. Đôi vai cô bé run lên từng chặp. Những ngón tay mũm mĩm thường ngày hay huơ lên đầy phấn khích nay bấu chặt vào gấu áo của Lamy như tìm kiếm một điểm tựa.

Giọng của em ấy nghẹn đắng khi micro mở, từng câu chữ đứt quãng: ``Bọn mình\dots\ đã chuẩn bị sẵn sàng mọi thứ để đón cậu ấy trở lại rồi mà\dots\ Bọn mình đã tập luyện cùng nhau\dots\ để có thể cùng làm buổi stream chung đầu tiên\dots''

Cô không thể nói hết câu. Tiếng sụt sùi vỡ vụn rơi vào khoảng không tĩnh lặng của phòng thu.

Polka và Botan ngồi ở vị trí trung tâm. Hai người họ gánh trên vai trọng trách nặng nề nhất: làm chỗ dựa để buổi stream không sụp đổ hoàn toàn. Pol-pol nở một nụ cười gượng gạo trước ống kính - nụ cười méo mó của một người nghệ sĩ tạp kỹ đang cố che giấu sự tan nát trong lòng. Hai bàn tay cô chắp chặt trước ngực, giọng nói cố giữ vẻ lưu loát nhưng đôi lúc lại run rẩy một nhịp chênh phô: ``Bọn mình biết\dots\ đây là sự thật nghiệt ngã mà bốn đứa phải học cách chấp nhận. Nhưng bọn mình vẫn sẽ tiếp tục cố gắng\dots\ vì cậu ấy, và vì tất cả những ai đã tin tưởng vào Gen 5.''

Botan là người giữ giọng bình tĩnh nhất, ít nhất là vẻ ngoài của nàng Sư tử trắng toát lên như vậy. Em hít một hơi thật sâu, ánh mắt kiên định nhìn thẳng vào màn hình, gãy gọn tóm tắt lại quyết định tốt nghiệp của người bạn cùng thế hệ vì lý do sức khỏe và tinh thần. Thế nhưng, đứng ngay phía sau lưng em, tôi nhìn thấy bờ vai Botan khẽ căng cứng. Mỗi lần nhấn phím chuyển màn hình, ngón tay em đều dừng lại một nhịp dài ngập ngừng.

Khi Lamy nghẹn ngào không thốt nên lời, tôi khẽ bước tới, không lọt vào khung hình, nhẹ nhàng đặt gói khăn giấy vào lòng bàn tay em rồi vỗ khẽ lên vai Nene. Một cái chạm vai trầm lặng của người làm quản lý.

Botan nghiêng đầu về phía tôi, ánh mắt như muốn nhờ cậy một điểm tựa: ``Kuon-senpai\dots\ anh có thể nói đôi lời với mọi người không ạ?''

Tôi bước lên một bước, đứng cạnh mép bàn, cất giọng trầm ấm vào micro: ``Chào các bạn khán giả. Tôi là Hikawa Kuon. Với tư cách là Tổng quản lý của nhánh \hll, tôi hiểu rõ sự trống vắng và hụt hẫng mà các thành viên ở đây cũng như tất cả các bạn đang phải trải qua. Mất đi một người đồng đội, một người bạn thân thiết khi chặng đường vừa mới bắt đầu là một nỗi đau không dễ gì nguôi ngoai. Nhưng tôi tin rằng, từ mất mát này, sự gắn kết giữa bốn thành viên còn lại sẽ càng thêm bền chặt. Công ty và cá nhân tôi sẽ luôn đồng hành, bảo vệ và hỗ trợ các em hết mình trên con đường phía trước. Xin các bạn hãy tiếp tục dõi theo và che chở cho các thành viên còn lại của Thế hệ thứ Năm.''

Buổi stream chỉ kéo dài vỏn vẹn hơn mười phút rồi khép lại.

Ngay khoảnh khắc tín hiệu phát sóng vừa tắt ngúm, lớp phòng ngự kiên cường cuối cùng của Botan cũng vỡ vụn. Nàng Sư tử gục đầu xuống bàn làm việc, hai tay ôm lấy trán. Căn phòng chìm trong tiếng nấc nghẹn ngào của bốn cô gái trẻ.

Hôm ấy, họ chia tay một cái tên. Một khoảng trống vĩnh viễn không thể lấp đầy trong đội hình xuất phát.

Nhưng họ không hề biết rằng, người bạn ấy chưa từng rời bỏ âm nhạc, và định mệnh vẫn âm thầm dệt nên một sợi dây liên kết để kéo họ trở về bên nhau.

\begin{center}
    \uline{Một ngày cuối tháng Tư, năm 2022.}
\end{center}

\pov{Momosuzu Nene}

``Cảm ơn quý khách! Chúc quý khách buổi tối vui vẻ ạ!''

Tiếng chuông cửa tiệm tạp hóa vang lên một tiếng leng keng giòn tan sau lưng tôi. Tôi vừa bước ra ngoài vỉa hè vừa thở phào một hơi dài, phả ra làn khói mỏng manh giữa tiết trời se lạnh cuối tháng Ba của khu phố của phường Anwa, nơi anh Kuon sinh sống.

Hai tay tôi xách hai túi đồ nặng trịch. Một bên là thịt bò thái mỏng, nấm kim châm, cải thảo và đậu phụ non để lát nữa cùng chú Ryujin nấu một nồi lẩu shabu-shabu hoành tráng. Bên còn lại - chẳng cần nói cũng biết - là cả một kho đồ ăn vặt linh tinh từ khoai tây chiên, bánh gạo cho tới mấy chai nước ngọt có ga mà cái bà Cáo Sa Mạc Polka vừa mới mè nheo nhắn tin bắt tôi mua bằng được.

Nghĩ đến cái mặt nhõng nhẽo của bà ấy là tôi lại thấy buồn cười: ``Nene-chi yêu quý ơi\~{} Mua thêm bim bim cay cho tôi đi mà\~{} Đói xỉu rồi nè\~{}''

Cứ như thể tôi là chân chạy việc của bà ấy không bằng ấy!

Cơ mà thôi, hôm nay là ngày vui đặc biệt. Cả bốn đứa NePoLaBo chúng tôi đã lên kế hoạch từ tuần trước, xin phép công ty để có một ngày nghỉ trọn vẹn kéo sang nhà anh Kuon chơi. Từ hồi sinh nhật anh ấy năm ngoái đến giờ, chúng tôi mới lại có dịp tụ tập đông đủ tại ``căn cứ'' này. Cả buổi chiều chơi game Mario Kart với nhau vui nổ trời, anh Kuon bình thường nghiêm túc là thế mà lúc cầm tay cầm lái xe cũng ghê gớm phết, mười trận thì thắng tới sáu, làm cả bọn cười đến sái cả quai hàm. Ryuuhou vừa ăn bim bim vừa hò reo cổ vũ ầm ĩ, bé Naomi thì ngồi trong lòng chị Suzuno vỗ tay khen anh trai rối rít, chỉ có chị Kaigou là khoanh tay lắc đầu trước cảnh tượng náo nhiệt ấy.

Tối nay chúng tôi còn dự định xin ngủ lại nữa. Anh Kuon đã gật đầu đồng ý, chú Ryujin và chị Suzuno thì hào phóng bảo sẽ đãi cả bọn một bữa lẩu nóng hổi cho ấm bụng.

Nghĩ đến nồi nước dùng sôi sùng sục thơm lừng mùi dashi, bước chân tôi bất giác thoăn thoắt hơn. Không biết chú Ryujin và chị Suzuno đã chuẩn bị xong phần nước chấm vừng mè chưa nữa--

``Ui da!''

Vì mải nghĩ lung tung và cắm đầu nhìn về phía trước, mũi giày của tôi va thẳng vào một gờ đá nhô cao trên vỉa hè.

Mất thăng bằng đột ngột, toàn bộ trọng tâm đổ nhào về phía trước. Tôi ngã sóng soài xuống nền xi măng, hai túi nilon trên tay văng ra xa. Tiếng túi nhựa rách toạc, những hộp nấm, củ cải trắng và gói snack khoai tây cứ thế lăn lông lốc khắp mặt đường vắng.

``Đau quá đi mất\dots''

Tôi rên rỉ một tiếng, nhăn nhó chống tay ngồi dậy. Cú ngã khiến hai đầu gối tôi nhói buốt, lớp vải quần tất rách một mảng nhỏ để lộ vết xước rớm máu. Bụi đất bám đầy vào lòng bàn tay và mép áo khoác.

Đúng là cái số hậu đậu muôn thuở của Nene này mà! Vừa đau, vừa tiếc túi đồ ăn, tôi chỉ muốn ngồi thụp xuống ăn vạ giữa đường cho bõ tức.

Thế nhưng, giữa lúc tôi còn đang xuýt xoa xoa nắn đầu gối, một bóng người từ phía trước bước nhanh tới, dừng lại ngay cạnh tôi.

Một bàn tay thanh mảnh, trắng trẻo chìa ra trước tầm mắt: ``Này, cậu có sao không? Đứng dậy được không đấy?''

Tôi giật bắn mình.

Không phải vì sự xuất hiện bất ngờ của người qua đường.

Mà là vì chất giọng ấy.

Một âm sắc vô cùng đặc biệt. Giọng nói mang âm vực cao trong trẻo nhưng ở cuối câu lại hơi khàn nhẹ, kèm theo một nhịp lấy hơi rất êm trước khi nhả chữ - thói quen của một người luyện thanh chuyên nghiệp. Đó là thứ âm thanh mà tôi đã khắc sâu vào trong tiềm thức suốt hai năm ròng rã, thứ giai điệu mà tôi ngỡ như sẽ chẳng bao giờ còn được nghe thấy ngoài đời thực nữa.

Tôi sững sờ ngẩng phắt đầu lên.

Đứng trước mặt tôi dưới ánh đèn đường vàng nhạt là một cô gái trẻ trạc tuổi tôi. Cô mặc một chiếc áo hoodie rộng thùng thình màu xám tro, quần jeans tối màu năng động. Mái tóc ngắn ngang vai mang sắc xanh lá mạ thẫm ánh khói, chiếc tai nghe chụp chuyên dụng vắt hờ hững lệch một bên cổ áo.

Không có cặp sừng quỷ màu đỏ. Không có chiếc đuôi nhọn tinh nghịch. Không có bộ váy succubus diêm dúa của thế giới ảo năm nào.

Người đứng trước mặt tôi là một cô gái hoàn toàn bằng xương bằng thịt, mang phong thái tự do, phóng khoáng của một ca sĩ đời thường.

Nhưng rồi, ánh mắt tôi trượt xuống bờ hông của cô.

Bên chiếc túi vải đeo chéo, một chiếc móc khóa nhỏ bằng kim loại đung đưa nhè nhẹ. Đó là hình một chiếc sừng quỷ cách điệu tinh xảo, với đường viền sơn phủ tông màu hồng pha tím quen thuộc. Bàn tay còn lại của cô, theo một thói quen vô thức khi đứng trước người lạ, đang ngượng ngùng đút sâu vào túi áo hoodie.

Thấy tôi cứ ngây người nhìn chằm chằm không chớp mắt, cô gái hơi nghiêng đầu sang một bên. Đôi mắt hồng tím long lanh sau làn tóc mái khẽ cong lên thành một nụ cười tinh quái, quen thuộc đến mức khiến tim tôi thắt nghẹn lại: ``Bà nhìn đi đâu thế hả? Nene-chi, ngã có một cú thôi mà đơ cả người ra rồi à?''

``Nene-chi.''

Tiếng cười khúc khích nhỏ nhẹ như chuông gió theo sau cái tên ấy như một mồi lửa châm thẳng vào ký ức.

Thời gian quanh tôi dường như đông cứng lại. Tiếng gió đêm, tiếng xe cộ xa xăm ngoài đại lộ hoàn toàn biến mất.

``M\dots\ Mano-chan\dots?''

Cổ họng tôi nghẹn ứ, hai môi run bần bật. Tôi chớp mắt liên tục, rồi bất giác giơ cả hai bàn tay lên, vỗ bôm bốp vào hai bên má mình. *BỐP! BỐP!*

Đau. Đau rát thật sự. Nhưng bóng hình trước mặt vẫn không hề tan biến.

``Kìa! Nene-chi, bà bị làm sao thế!?'' cô gái hốt hoảng cúi rạp người xuống, hai tay vội vàng giữ chặt lấy cổ tay tôi, ngăn không cho tôi tự tát mình nữa. ``Sao tự dưng lại tự đánh mình như thế chứ!?''

Làn da ấm áp chạm vào cổ tay tôi. Hơi thở gấp gáp của cô phả nhẹ qua gò má.

Đây không phải là một giấc mơ. Càng không phải là ảo ảnh do tôi nhớ bạn quá mà hóa thành.

Là cậu ấy. Bằng xương bằng thịt. Đang ở ngay trước mắt tôi.

Một luồng cảm xúc mãnh liệt như sóng thần trào dâng từ lồng ngực, đánh tan mọi sự e dè và hoài nghi. Tôi không thèm để ý đến vết xước ở đầu gối nữa, vùng mạnh khỏi tay cô ấy rồi chồm tới, vòng cả hai tay ôm chầm lấy cổ cô, siết chặt bằng tất cả sức lực của mình.

``Oái--!''

Cả hai chúng tôi cùng ngã ngửa ra nền đất. Nhưng tôi không buông ra. Tôi vùi mặt vào hõm vai thơm mùi xà phòng dịu nhẹ của cô, òa khóc nức nở như một đứa trẻ lạc vừa tìm thấy người thân: ``Oaaaa\dots! Mano-chan! Đúng là bà thật rồi! Hu hu hu\dots\ Bà có biết bọn tôi nhớ bà đến mức nào không hả!? Bà có biết bọn này đã lo lắng cho bà ra sao không\dots!''

Những giọt nước mắt ấm nóng tuôn trào không ngớt, thấm đẫm cả một mảng áo hoodie của cô bạn.

Người trong lòng tôi ban đầu cứng đờ người vì bất ngờ. Nhưng rồi, rất chậm rãi, tôi cảm nhận được hai cánh tay cô khẽ vòng qua lưng tôi, siết lại thật chặt. Đôi vai cô khẽ run rẩy theo từng nhịp thở ngắt quãng.

Một bàn tay nhẹ nhàng xoa lấy mái tóc tôi. Tiếng cô ấy thì thầm bên tai, vừa nghẹn ngào vừa chan chứa sự dịu dàng: ``Tôi biết rồi\dots\ Tôi xin lỗi nhé, Nene-chi\dots\ Tôi cũng nhớ mọi người nhiều lắm.''

\ct

Phải mất một hồi lâu sau, khi tiếng khóc của tôi đã vơi bớt thành những đợt sụt sùi ngắt quãng, hai đứa mới chịu buông nhau ra.

Mano - à không, giờ nhìn gần, phong thái của cậu ấy mang đậm nét của Delutaya hơn rất nhiều - lấy ngón tay gạt đi những giọt nước mắt còn vương trên khóe mắt tôi, rồi chìa tay đỡ tôi đứng dậy. Cậu ấy cúi xuống, nhặt từng hộp nấm và gói bánh vung vãi trên đất, phủi sạch bụi rồi xếp cẩn thận vào một chiếc túi mới.

``Đầu gối trầy da rồi kìa. Lát về phải sát trùng ngay đấy nhé,'' cậu ấy nhắc nhở, giọng nói đầy vẻ quan tâm.

``Tôi không sao mà!'' tôi quệt mũi, mắt vẫn không rời khỏi cậu ấy nửa bước, như thể chỉ cần chớp mắt một cái là người bạn này sẽ biến mất vậy. ``Cơ mà\dots\ sao bà lại ở đây? Khu này vắng lắm mà?''

Delutaya khẽ mỉm cười, đôi mắt lấp lánh ánh đèn đêm: ``Tôi vừa làm xong ca ở quán cà phê gần đây, đang đi dạo hóng gió một chút trước khi về phòng thu. Còn bà, ôm cả đống đồ ăn thế này đi đâu đấy?''

``Tôi đang đi mua đồ về nấu lẩu! Hôm nay cả đám NePoLaBo tụ tập bên nhà Kuon-senpai chơi nè! Mọi người đang đợi ở trên nhà rồi!''

Nói đến đây, một ý nghĩ lóe lên như tia chớp trong đầu tôi. Tôi chộp lấy cánh tay cậu ấy, mắt sáng rực: ``Này! Hay là bà sang chơi với bọn tôi luôn đi! Mọi người mà thấy bà chắc chắn sẽ mừng phát khóc cho xem!''

Nụ cười trên môi Delutaya bỗng chốc khựng lại. Cậu ấy hơi lùi lại nửa bước, ánh mắt thoáng vẻ bối rối, bàn tay vô thức lại đút sâu vào túi áo khoác: ``Ơ\dots\ Tôi á? Đến chỗ mọi người đang ở sao?''

``Ừ!''

``Nhưng mà\dots\ liệu có phiền không?'' cậu ấy ngập ngừng, hàng mi cụp xuống vẻ e dè. ``Đã lâu như vậy rồi\dots\ Mọi người liệu có thực sự muốn gặp một người đã rời đi như tôi không\dots?''

``Bà nói ngốc nghếch gì thế hả!'' tôi nắm chặt lấy tay cậu ấy, gằn từng chữ đầy quả quyết. ``Không một ai trong bọn tôi coi bà là người ngoài cả! Bọn tôi lúc nào cũng nhắc về bà, lúc nào cũng mong được gặp lại bà! Đi với tôi!''

Nhìn thấy sự kiên định không chút lung lay trong mắt tôi, sự do dự của Delutaya dần tan biến. Cậu ấy khẽ cắn môi, rồi nở một nụ cười ngượng ngùng: ``Được rồi\dots\ Nếu bà đã nói thế. Mà tiện thể\dots\ nhà của Kuon-kun mà bà nói, có phải là tòa nhà mặt phố viễn thông kia không?''

Cậu ấy vừa nói vừa giơ tay chỉ về phía con phố sáng đèn cách đó hai dãy nhà.

Tôi ngơ ngác chớp mắt: ``Ủa? Sao bà biết rành thế?''

Delutaya bật cười khúc khích, nhún vai một cái đầy vẻ trêu ngươi: ``Thì chúng ta đang đi chung một đường mà. Tôi đang sống ở tầng hai tòa nhà đó suốt hơn một năm nay rồi.''

Tôi đứng chết trân giữa vỉa hè, mồm há hốc ra đến mức có thể nhét vừa một quả trứng gà: ``\textbf{Hảaaaa!? Bà sống chung nhà với Kuon-senpai á!?}''

\ct

Sau một chuyến đi bộ ngắn, chúng tôi đã đứng trước cửa phòng tầng năm nhà Hikawa.

Tôi hít một hơi thật sâu, hai tay giữ chặt chiếc túi lẩu, quay sang nhìn người bạn đang đứng nép phía sau mình. Delutaya lúc này có vẻ hồi hộp thấy rõ. Hai bàn tay cậu ấy siết chặt lấy quai túi vải, ngón tay miết nhẹ vào lớp vải thô - một biểu hiện quen thuộc mỗi khi cậu ấy lo lắng trước giờ lên sóng.

``Bà sẵn sàng chưa?'' tôi thì thầm hỏi.

Delutaya nuốt khan một ngụm nước bọt, khẽ gật đầu: ``Ừ\dots\ Tôi sẵn sàng rồi.''

Tôi giơ tay đẩy mạnh cánh cửa trượt sang một bên, cất giọng oanh tạc quen thuộc vang dội khắp căn nhà: ``\textbf{Tôi về rồi đây cả nhà ơi! Đoán xem em mang ai về cùng này!}''

Tiếng bước chân từ trong phòng khách lập tức vang lên rộn rã. Bé Naomi đang ngồi vẽ tranh ở bàn trà liền ngẩng lên reo: ``A, chị Nene về rồi!''

``Nene-chi về rồi à! Có mua bim bim cay cho tôi không đấy-- Ủa?''

Polka là người chạy ra trước tiên, trên tay vẫn còn cầm một nửa chiếc bánh quế. Ngay sau lưng cô là Lamy đang bưng khay nước ngọt, Botan đang lững thững bước ra cùng Ryuuhou. Chị Suzuno cùng chú Ryujin đang bưng đĩa thịt bò từ bếp ra, còn chị Kaigou thì ngồi ở ghế sofa kiểm tra lại tập sổ sách.

Anh Kuon cũng từ trong bếp thò đầu ra, trên tay cầm chiếc muôi múc canh: ``Về đúng lúc đấy, nước lẩu vừa sôi--''

Lời nói của anh lửng lơ giữa chừng.

Cả căn phòng khách bỗng chốc rơi vào một sự tĩnh lặng tuyệt đối.

Delutaya từ từ bước ra từ sau lưng tôi, đứng dưới ánh đèn phòng khách ấm áp. Mái tóc xanh thẫm khẽ lay động, chiếc tai nghe vắt hờ một bên cổ áo, và chiếc móc khóa sừng quỷ lấp lánh phản chiếu ánh đèn.

Không ai trong ba người họ nhận ra ngay tức khắc. Trước mắt họ chỉ là một cô gái lạ mặt với trang phục đời thường, hoàn toàn không giống với hình tượng Mano Aloe trong ký ức năm nào.

Botan nheo mắt nhìn, khóe môi khẽ nhếch lên định buông lời trêu đùa: ``Gì đây Nene, đi siêu thị mà dắt cả bạn mới về chơi à--''

Nhưng cô chưa kịp nói hết câu.

Delutaya rụt rè giấu một tay vào túi áo hoodie, khẽ nghiêng đầu sang một bên. Cậu ấy hít một nhịp hơi rất nhẹ, rồi cất tiếng chào bằng chất giọng quen thuộc mà cả bốn chúng tôi từng cùng nhau luyện tập suốt bao ngày đêm: ``A-nô\dots\ Chào mọi người. Đã lâu không gặp\dots\ Botan-chan, Lamy-chan, Polka-chan.''

Khoảnh khắc âm thanh ấy vang lên, Lamy là người đầu tiên khựng lại.

Đôi tai nhạy cảm của cậu ấy khẽ giật giật. Chiếc khay đựng nước trên tay em chao đảo dữ dội. Là một người cực kỳ tinh tế với âm thanh, từng tần số, từng âm sắc trong trẻo pha chút khàn nhẹ của giọng nói ấy lập tức đập thẳng vào tâm trí Lamy. Đôi mắt màu vàng mở to hết cỡ, đồng tử co lại vì chấn động.

Tiếp theo đó là Polka. Trực giác nhạy bén của Cáo Sa Mạc mách bảo điều bất thường. Mắt cậu ấy nhìn trân trân vào thói quen giấu tay vào túi áo của cô gái, rồi lướt xuống chiếc móc khóa sừng quỷ viền tím hồng đung đưa bên hông. Miếng bánh quế trên tay cô tuột khỏi ngón tay, rơi bịch xuống sàn nhà.

``Kh\dots\ Không thể nào\dots'' giọng Polka run rẩy như sắp vỡ vụn.

Delutaya khẽ mỉm cười, đôi mắt hoe đỏ, giọng nói nghẹn ngào gọi lại lần nữa: ``Là tôi đây mà\dots\ Mano đây\dots''

*XOẢNG!*

Khay nước trên tay Lamy rơi thẳng xuống chiếc bàn thấp cạnh cửa. Không một ai bận tâm đến tiếng động ấy.

Chiếc đập ngăn cách cảm xúc suốt gần hai năm trời trong tích tắc vỡ tan tành.

``\textbf{Mano-chaaaan!}''

Lamy là người đầu tiên bùng nổ. Cậu ấy òa khóc nức nở như một đứa trẻ, lao thẳng tới không một chút do dự, ôm chầm lấy Delutaya với tất cả nỗi nhớ nhung dồn nén. Lực ôm mạnh đến mức khiến Delutaya phải lùi lại một bước để giữ thăng bằng.

``Hu hu hu\dots\ Đúng là cậu rồi\dots\ Cậu thật sự đã quay lại rồi\dots!'' Lamy vùi đầu vào vai cô gái, nước mắt tuôn rơi như mưa rào.

Polka cũng không thể kìm nén thêm được nữa. Cô hét lên một tiếng đầy nghẹn ngào rồi cũng nhào tới, vòng hai tay ôm ghì lấy cả Lamy lẫn Delutaya, vừa khóc vừa cười toe toét trong sự hỗn loạn: ``Con quỷ ngốc này! Đồ quỷ ngốc nhà cậu! Cậu trốn đi đâu suốt ngần ấy thời gian hả!? Cậu có biết bọn này nhớ cậu muốn phát điên lên không!''

Chỉ còn lại Botan.

Nàng Sư tử trắng đứng chôn chân ở ngưỡng cửa phòng khách. Đôi mắt màu xám của cậu ấy mở lớn, nhìn cảnh tượng ba người đang ôm nhau khóc lóc. Bờ vai khẽ run lên. Botan vội vàng quay mặt đi, đưa mu bàn tay áo lên che ngang mắt, cố gắng giấu đi những giọt nước mắt hiếm hoi đang lặng lẽ trào ra nơi khóe mi.

Trong ký ức của tôi, Botan gần như không bao giờ khóc trước mặt người khác. Cậu ấy luôn là chỗ dựa vững chãi nhất, là người mỉm cười động viên cả nhóm trong những thời khắc tăm tối nhất.

Nhìn thấy Botan như vậy, Delutaya nhẹ nhàng gỡ vòng ôm của Lamy và Polka ra. Cậu ấy bước từng bước chậm rãi đến trước mặt nàng Sư tử.

Delutaya ngước nhìn cậu ấy, hai bàn tay nắm chặt lấy gấu áo của chính mình, giọng nói thì thầm run rẩy: ``Botan-chan\dots\ Nếu tôi xuất hiện ở đây làm bà thấy phiền lòng\dots''

Không để cậu ấy nói hết câu, Botan đã đột ngột quay người lại. Bằng một động tác dứt khoát nhưng tràn đầy tình cảm, cô dang rộng vòng tay, kéo phắt Delutaya vào một cái ôm chặt cứng vào lồng ngực mình.

Bàn tay to lớn của nàng Sư tử xoa mạnh lên mái tóc ngắn màu xanh của cô em gái, tiếng cười khẽ xen lẫn tiếng nấc nghẹn ngào: ``Nói ngốc xít gì thế hả con bé này\dots\ Phiền cái gì chứ\dots\ Mừng cậu trở về nhà, Mano\dots\ à không, Delutaya.''

Bốn đứa chúng tôi lại một lần nữa lao vào, ôm chầm lấy nhau thành một vòng tròn ấm áp ngay giữa phòng khách tầng năm. Tiếng khóc, tiếng cười, tiếng gọi tên nhau vang vọng khắp gian phòng, xóa nhòa đi mọi khoảng cách thời gian và không gian của hai năm chia cách.

Bé Naomi - xuất hiện từ đâu đó mà chúng tôi không để ý - tròn xoe mắt reo lên: ``A! Delu-oneechan ở tầng hai lên kìa!''

Suzuno hai tay ôm khay chén đĩa, mắt rưng rưng mỉm cười dịu dàng. Chị Kaigou gấp quyển sổ lại, ánh mắt nghiêm nghị thường ngày khẽ dịu đi, khóe môi khẽ cong lên một nụ cười nhẹ nhõm. Ryuuhou tháo phắt một bên tai nghe xuống, hào sảng vỗ tay: ``Hay lắm! Hôm nay khách quý tầng hai cũng lên ăn tối cùng cả nhà rồi!''

Đứng ở cửa bếp, anh Kuon mỉm cười hiền từ, khẽ gật đầu với bố Ryujin. Chú Ryujin cũng bật cười sảng khoái, vỗ đùi đánh đét một cái: ``Tốt rồi! Tốt lắm! Khóc lóc thế đủ rồi mấy đứa, mau vào bàn ăn lẩu kẻo nước cạn hết bây giờ!''

\ct

Nồi lẩu shabu-shabu bốc khói nghi ngút ở giữa bàn ăn, tỏa ra mùi thơm ngọt ngào của nước dùng xương hầm cùng rau củ tươi.

Dưới tiết trời se lạnh cuối xuân, được quây quần bên nồi lẩu nóng hổi thế này quả thực không còn gì ấm áp bằng. Cả nhà Hikawa cùng năm cô gái ngồi quây quần kín quanh chiếc bàn gỗ lớn kéo dài. Bé Naomi ngồi cạnh chị Suzuno, được chị gắp cho phần nấm và viên chả tôm thơm phức; Ryuuhou và Kaigou thỉnh thoảng lại tranh nhau đĩa đậu phụ non.

Lamy ngồi ngay sát cạnh Delutaya, hết gắp miếng thịt bò ngon nhất cho cậu ấy lại rót đầy ly nước ấm. Cô bạn Tam giác cầm lấy ly nước, ngón tay trỏ lại miết nhẹ theo đường viền miệng ly - cử chỉ quen thuộc mỗi khi cậu ấy ngượng ngùng trước sự quan tâm quá mức của người khác.

``Nào, uống chút nước ấm đi Mano\dots\ à nhầm, Delu-chan,'' Lamy dịu dàng nói.

``Cảm ơn Lamy-chan\dots\ Mà mọi người cứ gọi tôi là Delutaya hay Aloe gì cũng được mà, đâu cần phải câu nệ thế,'' cậu ấy mỉm cười đáp lại.

Polka vừa nhai tóp tép miếng thịt bò vừa chĩa đôi đũa về phía Kuon-senpai, vẻ mặt đầy sự chất vấn: ``Cơ mà này Kuon-senpai! Anh giải thích rõ ràng cho bọn em coi! Tại sao Aloe\dots\ à nhầm, Delutaya lại ở ngay tầng hai nhà anh suốt thời gian qua mà anh không hé răng nửa lời với bọn em hả!?''

Botan cũng phụ họa, nheo mắt nhìn anh quản lý: ``Đúng đó Kuon-senpai. Bọn em đến nhà anh chơi bao nhiêu lần, thế mà anh giấu kỹ quá đấy nhé.''

Anh Tổng bất đắc dĩ đặt bát mì xuống, xua xua tay: ``Oan cho anh quá mấy đứa ơi. Kế hoạch này là do Yagoo-san và A-san sắp xếp từ trước. Nơi này được chọn làm một dạng `ký túc xá an toàn' để các cựu thành viên sau khi tốt nghiệp có thể an tâm sinh sống và tiếp tục theo đuổi đam mê mà không bị quấy rầy. Chú ấy muốn mọi chuyện diễn ra tự nhiên nhất có thể, khi tinh thần của Delutaya đã dần ổn định hơn và sẵn sàng đối mặt. Bản thân anh cũng đâu ngờ Nene lại tình cờ gặp em ấy giữa đường như thế.''

Chú Ryujin húp một ngụm canh, cười ha hả: ``Thằng Kuon nó nói thật đấy các cháu. Tiền phòng hàng tháng của Delutaya đều do ông ấy đích thân chuyển khoản đàng hoàng, không thiếu một xu. Chú thấy con bé nó chăm chỉ lắm, ngày nào cũng luyện giọng với làm việc ở quán cà phê Domino đấy.''

``Quán Domino ngay cạnh đây sao!?'' tôi ngạc nhiên kêu lên. ``Thảo nào đợt trước em thấy bóng ai quen quen đứng ở quầy pha chế!''

Delutaya gật đầu, nụ cười trên môi đã trở nên tự nhiên và rạng rỡ hơn bao giờ hết: ``Ừ, tôi làm phục vụ ở đó để trang trải thêm sinh hoạt phí. Ngoài ra thì\dots\ như Kuon-senpai và Ryujin-san đã biết, tôi vẫn tiếp tục hoạt động với tư cách VTuber độc lập mang tên Delutaya.''

``Bọn này biết chứ!'' Polka phấn khích đập bàn. ``Kênh của bà bọn tôi đều âm thầm theo dõi hết đấy nhé! Những buổi stream ca hát của bà tuyệt vời lắm!''

``Thật sao\dots?'' cô bạn Tam giác chớp mắt, hai má ửng hồng.

``Đương nhiên rồi!'' Botan mỉm cười ấm áp. ``Nghe nói cậu sắp chuẩn bị ra mắt mô hình 3D mới đúng không?''

``Vâng. Đợt vừa rồi tôi nhờ Kuon-senpai và mọi người ở đây hỗ trợ dựng sân khấu chuyển động để tổng duyệt\footnote{Chi tiết trong {\flaregothic HISTOIRE \freshbot 19} thuộc quyển này.}. Mọi thứ đã gần như hoàn tất rồi. Sắp tới tôi sẽ chính thức tái xuất.''

Tôi nắm lấy bàn tay cậu ấy lắc mạnh, hào hứng reo lên: ``Nhất định hôm đó bọn tôi sẽ xem! Bọn tôi sẽ rủ cả fan vào cày view cho bà luôn!''

``Cảm ơn mọi người nhiều lắm\dots''

Delutaya cúi đầu nhìn làn khói nghi ngút bốc lên từ nồi lẩu, ánh mắt trầm lắng lại trong giây lát. Cậu ấy khẽ thở ra một hơi dài, rồi cất lời bằng tất cả sự chân thành: ``Mọi người biết không\dots\ Vào cái ngày cuối cùng ấy, ngày ba mươi mốt tháng Tám năm 2020\dots\ khi xem buổi stream thông báo của bốn người, tôi đã khóc rất nhiều. Tôi thấy Lamy-chan và Nene-chi rơi nước mắt vì tôi, thấy Polka-chan và Botan-chan cố gượng cười để giữ vững tinh thần cho cả nhóm\dots\ Chính tình cảm và sự kiên cường của các cậu ngày hôm đó đã trở thành động lực lớn nhất giúp tôi không bỏ cuộc. Nhờ có bốn người, tôi mới có dũng khí để tiếp tục đứng dậy và cất tiếng hát cho đến tận ngày hôm nay.''

Tôi bỗng thấy cay cay nơi khóe mắt, liền chồm tới phản bác: ``Bà đừng có nói thế! Bà cũng đã rất kiên cường chứ bộ! Tự mình vượt qua bao nhiêu khó khăn, bao nhiêu áp lực từ những chuyện không hay ngày ấy, dũng cảm tiếp tục bước đi với đam mê âm nhạc, bà siêu lắm rồi đấy!''

Polka cũng bừng bừng khí thế, gằn từng chữ: ``Đúng vậy! Cứ nghĩ đến những khó khăn và áp lực mà bà phải chịu đựng hồi đó là tôi lại tức anh ách! Mấy kẻ chỉ biết ngồi sau bàn phím buông lời cay nghiệt gây khó dễ\dots''

Botan siết chặt nắm tay, trong mắt lóe lên ngọn lửa bảo bọc quyết liệt: ``Nếu sau này còn có kẻ nào dám gây khó dễ hay bắt nạt Delu-chan nữa, chị em chúng ta nhất định sẽ không để yên đâu! Phải cho chúng biết bọn này không dễ bị bắt nạt!''

Lamy cũng quẹt nước mắt, gật đầu lia lịa: ``Đúng thế! Bọn mình sẽ bảo vệ cậu bằng mọi giá!''

Cả bốn cô nàng NePoLaBo chúng tôi hăng máu ngút trời, đồng loạt đứng phắt dậy, vung tay định đập bàn hô vang lời thề bảo vệ chị em--

*RẦM!*

Vì đứng dậy quá đột ngột và hăng máu, đầu gối của tôi và Polka va thẳng vào mép bàn gỗ nặng trịch! Chiếc bàn rung chuyển dữ dội, chiếc kiềng bếp ga mini trượt lệch hẳn sang một bên, và nồi lẩu shabu-shabu đang sôi ùng ục chao đảo nghiêng ngả, nước dùng nóng hổi suýt chút nữa là đổ ụp xuống sàn nhà.

``\textbf{Ối giời đất ơi!!}''

Cả căn phòng khách rơi vào tình trạng báo động đỏ.

Anh Kuon và chị Kaigou phản xạ như chớp, mỗi người một tay lao tới tóm chặt lấy hai quai nồi lẩu và ghì chặt mép bàn.

Ryuuhou giật bắn mình suýt rơi cả đũa, vội giơ hai tay đỡ lấy chiếc bếp ga: ``Này này này! Mấy chị bình tĩnh lại giùm con với! Suýt nữa là luộc chín cả phòng khách rồi đấy!''

Suzuno nhanh thoăn thoắt kéo đĩa thịt bò và bát nước chấm giật lùi về phía sau, bảo vệ đồ ăn trong gang tấc.

Bé Naomi hoảng hốt ôm chặt lấy cánh tay chị mình, mắt tròn xoe nhìn chiếc nồi vẫn còn bốc khói nghi ngút.

Chú Ryujin thì ôm ngực toát mồ hôi hột, thở hắt ra một hơi: ``Hú hồn con chim én! Chú biết các cháu có tinh thần hiệp nghĩa bảo vệ bạn bè rồi, nhưng làm ơn để cho cái nồi lẩu nó sống sót qua tối nay đã! Cháy cái thảm này là chú bắt thằng Kuon đền đấy nhé!''

Anh Kuon dở khóc dở cười, tay vẫn còn giữ chặt quai nồi: ``Anh rất cảm kích nhiệt huyết của mấy đứa, nhưng trước khi đòi đi cho ai một bài học thì làm ơn giữ mạng cho cái nồi lẩu tối nay trước đã\dots''

Bốn cô nàng NePoLaBo chúng tôi đang đứng bỗng sượng sùng cứng đờ người. Nhìn nồi lẩu suýt tan tành vì sự quá khích của mình, cả bốn đứa - cùng Delutaya đang ngồi bên cạnh - ngượng chín cả mặt.

Năm đứa nhìn nhau, rồi lí nhí cúi đầu đồng thanh: ``D-Dạ\dots\ bọn cháu xin lỗi cả nhà ạ\dots''

Bé Naomi thấy cảnh các chị cúi đầu hối lỗi thì không nhịn được, bật cười khúc khích. Tiếng cười trong trẻo của cô bé xua tan đi cơn hoảng loạn, kéo theo tràng cười sảng khoái của cả phòng khách.

Botan hắng giọng một tiếng xua đi sự ngượng ngùng, rồi với tay qua bàn, đặt bàn tay ấm áp của mình lên mu bàn tay của Delutaya: ``Đó là vì chúng ta là một gia đình mà. Dù cậu mang tên gọi nào, dù hoạt động ở công ty hay độc lập bên ngoài, đối với chúng tôi\dots\ cậu mãi mãi là một mảnh ghép không thể thiếu.''

``Đúng vậy!'' Lamy quẹt nước mắt, mỉm cười rạng rỡ. ``holoFive chúng ta mãi mãi là năm người!''

``Nâng ly nào chị em ơi!'' Polka giơ cao cốc nước ngọt.

Những chiếc ly cùng cụng vào nhau chan chát giữa tiếng cười giòn giã. Hơi ấm của nồi lẩu, của tình bạn và sự sẻ chia lan tỏa khắp căn phòng, xua tan đi cái lạnh lẽo của màn đêm ngoài khung cửa sổ.

Một buổi tái ngộ trọn vẹn và đẹp đẽ hơn bất cứ điều gì chúng tôi từng dám mơ ước.

%==========================================TRUYỆN PHỤ=========================================
\omk

\pov{-}

Bữa tiệc lẩu kết thúc khi đồng hồ điểm sang chín giờ tối.

Sau khi đã phụ giúp chị Suzuno và chú Ryujin dọn dẹp sạch sẽ bát đĩa và nồi niêu trên bếp, bốn cô gái NePoLaBo liền tụ tập lại ở phòng khách, vẻ mặt ai nấy đều lộ rõ sự tiếc nuối khi phải nghĩ đến chuyện ra về.

Kuon từ trong phòng làm việc bước ra, nhìn bốn cô gái đang ngồi túm tụm: ``Cũng muộn rồi đấy. Mấy đứa có cần anh gọi xe taxi đưa về không?''

Nene lập tức phồng má, lắc đầu nguây nguẩy: ``Ể\~{} Bọn em không về đâu! Đã nói từ trước là tối nay bọn em ngủ lại đây mà!''

``Hả? Nhưng nhà anh làm gì còn phòng ngủ trống cho cả bốn đứa?'' Kuon nhíu mày.

Polka cười toe toét, chỉ ngón tay xuống sàn nhà: ``Thì xuống phòng Delu-chan ở tầng hai ngủ chứ đâu nữa anh! Phòng của bà ấy rộng thế cơ mà!''

Delutaya đang đứng bên cạnh bỗng giật nảy mình, hai tai đỏ ửng: ``Ơ\dots\ phòng của tôi sao!? Nhưng phòng tôi chưa dọn dẹp kỹ lắm\dots\ với lại chỉ có đệm trải sàn thôi\dots''

``Có đệm là tốt lắm rồi! Năm chị em mình ngủ chung một tấm đệm mới vui chứ!'' Nene reo lên đầy phấn khích, không để cho người bạn nhỏ có cơ hội từ chối.

Lamy quay sang nhìn Kuon, chớp chớp mắt đầy vẻ tinh nghịch: ``Kuon-senpai không định xuống ngủ cùng bọn em luôn sao?''

Mặt cậu Tổng lập tức đanh lại, anh lùi phắt lại một bước như vừa nhìn thấy quái vật: ``Mấy đứa bị dở hơi à!? Nam nữ thụ thụ bất thân, anh xuống đó làm cái gì!?''

``Gì chứ, sinh nhật năm ngoái anh chẳng ngủ chung phòng với tụi em là gì\~{}'' Polka phụ họa, cười khúc khích.

Kuon ôm trán rên rỉ bất lực trước sự trêu chọc của mấy cô nàng tinh quái. Ryujin đứng dựa cửa bếp cười ha hả giải vây cho con trai: ``Thôi thôi, tha cho thằng Kuon đi các cháu. Để cho nó còn làm việc. Mấy đứa cứ thoải mái kéo nhau xuống tầng hai mà tâm sự tuổi hồng.''

``Cảm ơn chú Ryujin ạ! Đi thôi chị em ơi!''

Cả năm cô gái ríu rít như bầy chim non kéo nhau chạy ào xuống cầu thang tầng hai. Chị Suzuno mỉm cười ôm thêm hai chiếc chăn bông thơm tho và gối sạch mang xuống cho năm đứa, còn bé Naomi đứng ở bậc thang vẫy tay toe toét: ``Chúc chị Delu-chan với các chị ngủ ngon ạ!''

Chị Kaigou đứng khoanh tay cạnh anh Kuon, nghiêm giọng nhắc nhở với xuống: ``Chơi game vừa phải thôi đấy, mười một giờ là phải tắt đèn đi ngủ, ngày mai còn phải thức dậy đúng giờ.''

Ryuuhou tựa lưng vào tường, nhếch môi cười trêu: ``Kaigou-nee lại bắt đầu phát huy bệnh quản gia rồi đấy.''

``Chị đang nói điều hoàn toàn hợp lý, Ryuuhou,'' chị Kaigou lườm cô em gái một cái sắc lẹm, nhưng khóe môi cũng không giấu được nụ cười nhẹ nhõm.

\ct

Căn phòng tầng hai của Delutaya sáng đèn rực rỡ.

Không gian phòng không quá cầu kỳ, một góc là bàn làm việc với dàn máy tính cấu hình cao và micro thu âm chuyên nghiệp, góc còn lại là một chiếc đệm lớn trải sẵn trên sàn gỗ.

Botan nhanh thoăn thoắt cắm chiếc máy chơi game cầm tay vào màn hình phụ của Delutaya. Polka thì bày la liệt những túi bánh khoai tây chiên và nước ngọt ra giữa sàn, trong khi Nene và Lamy lăng xăng giúp cô chủ nhà trải thêm chăn gối.

Đêm hôm ấy, tiếng bấm nút lách cách, tiếng hò reo cổ vũ trong những trận đấu game kịch tính xen lẫn tiếng cười đùa vang vọng khắp căn phòng nhỏ. Họ chia nhau từng miếng snack, kể cho nhau nghe những câu chuyện vụn vặt suốt hai năm qua - những trò đùa ngớ ngẩn trên stream, những lo toan trong cuộc sống thường nhật, và cả những dự định rực rỡ cho tương lai.

Đến nửa đêm, khi những khay bánh đã cạn và mí mắt ai nấy đều trĩu nặng, cả năm người nằm chen chúc nhau trên chiếc đệm êm ái.

Năm bàn tay đan chặt vào nhau dưới lớp chăn ấm.

Không còn ranh giới giữa cựu thành viên và thành viên đương nhiệm. Không còn khoảng cách giữa ánh đèn sân khấu và cuộc sống đời thường.

Dưới mái nhà Hikawa đêm ấy, Thế hệ thứ Năm của \textit{hololive} đã thực sự tìm lại được nhau - trọn vẹn, ấm áp và vĩnh cửu.

\ct

\pov{Hikawa Kuon}

Mười một giờ đêm.

Tôi đứng ngoài ban công tầng năm, tựa tay vào lan can kim loại đón những cơn gió lạnh thổi qua rặng cây ven đường. Thành phố Kawachi về đêm chìm trong những dải ánh sáng vàng cam tĩnh lặng.

Tiếng chuông điện thoại trong túi áo rung lên nhè nhẹ.

Trên màn hình hiện lên hai chữ quen thuộc: YAGOO.

Tôi bấm nút nhận cuộc gọi, đưa máy lên tai: ``Hikawa Kuon nghe đây ạ.''

Đầu dây bên kia im lặng trong một giây, rồi giọng nói điềm đạm quen thuộc của người đàn ông đứng đầu công ty cất lên: ``Tình hình bên đó thế nào rồi, Kuon-kun?''

Tôi nhìn xuống ánh đèn hắt ra từ khung cửa sổ tầng hai - nơi vẫn còn văng vẳng tiếng cười khúc khích nhỏ nhẹ của năm cô gái.

``Báo cáo chú, chiến dịch Réunion bước đầu tiên đã thành công tốt đẹp. Bọn họ đã gặp lại nhau, cùng ăn một bữa tối trọn vẹn và hiện đang ngủ chung ở phòng Delu-chan rồi ạ.''

Một tiếng thở phào nhẹ nhõm vang lên từ đầu dây bên kia, theo sau là tiếng cười khẽ đầy ấm áp của sếp tôi: ``Tốt lắm. Cậu đã vất vả rồi. Tôi biết để bọn họ tự tìm thấy nhau là cách tốt nhất.''

``Vâng. Cảm xúc rất chân thật, không một chút gượng ép nào ạ. Mỗi tội\dots\ mọi thứ lại diễn ra quá bất ngờ. Cháu cũng không thể nghĩ được là mọi người sẽ tụ họp nhau theo cách đó ạ.''

Chú Tanigo bật cười khúc khích từ đầu dây, tiếng cười chất chứa niềm vui dồn nén bao lâu nay: ``Ngẫu nhiên mới là điều đẹp nhất phải không nào? Nếu chúng ta sắp xếp hết mọi chuyện theo kịch bản, liệu có được khoảnh khắc chân thành như đêm nay không?''

Tôi im lặng một lát, nhìn ánh đèn vàng ấm áp thoát ra từ khe cửa sổ tầng hai, rồi đáp nhẹ nhàng: ``Chú nói đúng. Nước mắt và nụ cười của họ đêm nay\dots\ là món quà quý giá nhất mà chúng ta có thể nhận được.''

``Vậy thì\dots'' giọng chú Tanigo trầm xuống một nhịp, mang theo sự suy tư sâu xa. ``Bước tiếp theo có lẽ chúng ta cũng nên bắt đầu chuẩn bị rồi nhỉ? Coco-chan\dots\ và Thế hệ thứ Tư.''

Tôi khẽ gật đầu, dù biết ông ấy không nhìn thấy: ``Kson-chan thì cháu nghĩ không có vấn đề gì lớn đâu. Em ấy đã rất vững vàng rồi.''

``Ừ. Nhưng còn Rushia-chan thì sao?''

Câu hỏi ấy khiến lòng tôi bất giác chùng xuống.

Tôi ngước mắt nhìn lên tầng tám của tòa nhà - nơi căn phòng của Mikeneko vẫn chìm trong bóng tối im lìm. Một vết thương quá sâu, một nỗi đau vẫn chưa thể lành miệng.

``Trường hợp của em ấy\dots\ sẽ cần thêm rất nhiều thời gian và kiên nhẫn ạ. Em ấy vẫn chưa thực sự sẵn sàng đâu,'' tôi đáp chậm rãi.

``Tôi hiểu. Cứ từng bước một thôi, Kuon-kun. Đêm nay, hãy cứ để các cô gái tận hưởng niềm vui trọn vẹn này đã.''

``Vâng ạ. Chúc chú ngủ ngon.''

Cuộc gọi kết thúc.

Tôi cất điện thoại vào túi áo, hít một hơi sâu luồng không khí thanh sạch của đêm xuân.

Ngay ở tầng hai bên dưới, bóng đèn huỳnh quang cuối cùng trong căn phòng nhỏ cũng nhẹ nhàng tắt lịm, để lại chỉ có ánh trăng bạc từ cửa sổ chiếu rọi lên năm gương mặt đang yên bình chìm vào giấc ngủ. Những tiếng cười đùa cuối cùng đã nhường chỗ cho hơi thở đều đặn, những tâm hồn từng trải qua biết bao sóng gió, chia ly và đơn độc, giờ đây cuối cùng cũng tìm được sự an toàn để neo đậu bên nhau sau bao lâu xa cách.

Đêm nay chỉ là trang đầu tiên của một chương mới tuyệt đẹp, nơi những vết thương cũ trong quá khứ sẽ dần được chữa lành, và những ước mơ bị bỏ dở một thời gian sẽ lại được tiếp tục viết nên bằng niềm vui trọn vẹn.

\end{document}
